% GENERATED FILE. Edit canonical lessons, then run npm run generate:books.
% canonical-source-hash: fnv1a64:c3aa1cb5cb116e80
% canonical-lessons: PT-C530-senhorio, PT-C530-inquilino, PT-C530-mobilar, PT-C530-despensa, PT-C530-humidade, PT-R530-revisao-523-530

\chapter{Senhorios e inquilinos}
\label{ch:pt-senhorios-e-inquilinos}
\hlchaptermodality{530}

\begin{chapteropening}
\textbf{By the end of this chapter:} \emph{I can talk about renting, furnishing and storage.}

Complete the last lesson of chapter 530: I can talk about renting, furnishing and storage.
\end{chapteropening}

\section[o senhorio]{o senhorio — landlord}

\label{lesson:PT-C530-senhorio}



\begin{quote}
Before the new one: say the Portuguese for clothes peg, then the Portuguese for to iron.
\end{quote}

\subsection*{You'll want to know: o senhorio}
\textbf{o senhorio} — \textquotedblleft{}landlord\textquotedblright{}.

\begin{quote}
O senhorio aumentou a renda.
\end{quote}

\begin{quote}
\emph{The landlord put up the rent.}
\end{quote}

\emph{A senhoria} — landlady.

\begin{cousinweb}
\emph{senhorio} is inherited from Old Portuguese \emph{sennorio} \textquotedblleft{}lordship\textquotedblright{} (EtymDB), built on \emph{senhor}, from Latin \emph{senior} \textquotedblleft{}older\textquotedblright{}.
\end{cousinweb}

\begin{grammarlens}[title={What you've built}]
One more word for telling, explaining and describing at B1.
\end{grammarlens}

\subsection*{Your turn}
\begin{itemize}
\raggedright
  \item \emph{Say it:} \emph{o senhorio}
  \item \emph{Say it:} \emph{o senhorio}, once more
  \item \emph{Say it:} a sentence of your own with \emph{o senhorio}
  \item \emph{Recall:} say the Portuguese for clothes peg, then the Portuguese for to iron, then say \emph{o senhorio} again
\end{itemize}

\begin{tcolorbox}[breakable,colback=teal!4,colframe=teal!35!black,title={Before you move on}]
What does o senhorio mean? (\textquotedblleft{}Landlord\textquotedblright{}.) Say it once more.
\end{tcolorbox}
\section[o inquilino]{o inquilino — tenant}

\label{lesson:PT-C530-inquilino}



\begin{quote}
Before the new one: say the Portuguese for to iron, then the Portuguese for landlord.
\end{quote}

\subsection*{You'll want to know: o inquilino}
\textbf{o inquilino} — \textquotedblleft{}tenant\textquotedblright{}.

\begin{quote}
Os inquilinos do segundo andar são muito simpáticos.
\end{quote}

\begin{quote}
\emph{The tenants on the second floor are very friendly.}
\end{quote}

A woman is \emph{a inquilina}.

\begin{cousinweb}
\emph{inquilino} is a learned borrowing from Latin \emph{inquilīnus} \textquotedblleft{}a tenant, a lodger\textquotedblright{} (EtymDB).
\end{cousinweb}

\begin{grammarlens}[title={What you've built}]
One more word for telling, explaining and describing at B1.
\end{grammarlens}

\subsection*{Your turn}
\begin{itemize}
\raggedright
  \item \emph{Say it:} \emph{o inquilino}
  \item \emph{Say it:} \emph{o inquilino}, once more
  \item \emph{Say it:} a sentence of your own with \emph{o inquilino}
  \item \emph{Recall:} say the Portuguese for to iron, then the Portuguese for landlord, then say \emph{o inquilino} again
\end{itemize}

\begin{tcolorbox}[breakable,colback=teal!4,colframe=teal!35!black,title={Before you move on}]
What does o inquilino mean? (\textquotedblleft{}Tenant\textquotedblright{}.) Say it once more.
\end{tcolorbox}
\section[mobilar]{mobilar — to furnish}

\label{lesson:PT-C530-mobilar}



\begin{quote}
Before the new one: say the Portuguese for landlord, then the Portuguese for tenant.
\end{quote}

\subsection*{You'll want to know: mobilar}
\textbf{mobilar} — \textquotedblleft{}to furnish\textquotedblright{}.

\begin{quote}
Mobilámos a casa com móveis usados.
\end{quote}

\begin{quote}
\emph{We furnished the house with second-hand furniture.}
\end{quote}

In Brazil: \emph{mobiliar}.

\begin{cousinweb}
\emph{mobilar} is built on \emph{mobília} \textquotedblleft{}furniture\textquotedblright{}, from Latin \emph{mōbilis} \textquotedblleft{}movable\textquotedblright{} (Lewis \& Short), from \emph{movēre} \textquotedblleft{}to move\textquotedblright{}.
\end{cousinweb}

\begin{grammarlens}[title={What you've built}]
One more word for telling, explaining and describing at B1.
\end{grammarlens}

\subsection*{Your turn}
\begin{itemize}
\raggedright
  \item \emph{Say it:} \emph{mobilar}
  \item \emph{Say it:} \emph{mobilar}, once more
  \item \emph{Say it:} a sentence of your own with \emph{mobilar}
  \item \emph{Recall:} say the Portuguese for landlord, then the Portuguese for tenant, then say \emph{mobilar} again
\end{itemize}

\begin{tcolorbox}[breakable,colback=teal!4,colframe=teal!35!black,title={Before you move on}]
What does mobilar mean? (\textquotedblleft{}To furnish\textquotedblright{}.) Say it once more.
\end{tcolorbox}
\section[a despensa]{a despensa — pantry, larder}

\label{lesson:PT-C530-despensa}



\begin{quote}
Before the new one: say the Portuguese for tenant, then the Portuguese for to furnish.
\end{quote}

\subsection*{You'll want to know: a despensa}
\textbf{a despensa} — \textquotedblleft{}pantry, larder\textquotedblright{}.

\begin{quote}
Guardamos o arroz e as latas na despensa.
\end{quote}

\begin{quote}
\emph{We keep the rice and the tins in the pantry.}
\end{quote}

Not to be mixed up with \emph{a dispensa}, an exemption.

\begin{cousinweb}
\emph{despensa} is usually traced to Medieval Latin \emph{dispēnsa} \textquotedblleft{}a store of provisions\textquotedblright{}, from \emph{dispendere} \textquotedblleft{}to weigh out\textquotedblright{}, built on \emph{pendere} \textquotedblleft{}to weigh\textquotedblright{} (Lewis \& Short).
\end{cousinweb}

\begin{grammarlens}[title={What you've built}]
One more word for telling, explaining and describing at B1.
\end{grammarlens}

\subsection*{Your turn}
\begin{itemize}
\raggedright
  \item \emph{Say it:} \emph{a despensa}
  \item \emph{Say it:} \emph{a despensa}, once more
  \item \emph{Say it:} a sentence of your own with \emph{a despensa}
  \item \emph{Recall:} say the Portuguese for tenant, then the Portuguese for to furnish, then say \emph{a despensa} again
\end{itemize}

\begin{tcolorbox}[breakable,colback=teal!4,colframe=teal!35!black,title={Before you move on}]
What does a despensa mean? (\textquotedblleft{}Pantry, larder\textquotedblright{}.) Say it once more.
\end{tcolorbox}
\section[a humidade]{a humidade — damp, humidity}

\label{lesson:PT-C530-humidade}



\begin{quote}
Before the new one: say the Portuguese for to furnish, then the Portuguese for pantry.
\end{quote}

\subsection*{You'll want to know: a humidade}
\textbf{a humidade} — \textquotedblleft{}damp, humidity\textquotedblright{}.

\begin{quote}
Há humidade nas paredes do quarto.
\end{quote}

\begin{quote}
\emph{There is damp on the bedroom walls.}
\end{quote}

In Brazil: \emph{a umidade}.

\begin{cousinweb}
\emph{humidade} is built on \emph{húmido} \textquotedblleft{}damp\textquotedblright{}, a learned borrowing from Latin \emph{ūmidus} \textquotedblleft{}moist\textquotedblright{} (Lewis \& Short).
\end{cousinweb}

\begin{grammarlens}[title={What you've built}]
One more word for telling, explaining and describing at B1.
\end{grammarlens}

\subsection*{Your turn}
\begin{itemize}
\raggedright
  \item \emph{Say it:} \emph{a humidade}
  \item \emph{Say it:} \emph{a humidade}, once more
  \item \emph{Say it:} a sentence of your own with \emph{a humidade}
  \item \emph{Recall:} say the Portuguese for to furnish, then the Portuguese for pantry, then say \emph{a humidade} again
\end{itemize}

\begin{tcolorbox}[breakable,colback=teal!4,colframe=teal!35!black,title={Before you move on}]
What does a humidade mean? (\textquotedblleft{}Damp, humidity\textquotedblright{}.) Say it once more.
\end{tcolorbox}
\section[(dialogue)]{Revisão dos capítulos 523–530}

\label{lesson:PT-R530-revisao-523-530}



\begin{quote}
Say the Portuguese for chapter and novelist.
\end{quote}

\subsection*{Your turn}
Say these aloud:

\begin{itemize}
\raggedright
  \item chapter, novelist, biography, to reread, bookmark
  \item poetry, poet, line, rhyme, to recite
  \item librarian, to borrow, publisher, translation, translator
  \item episode, season, soap opera, subtitled, dubbing
  \item (film) director, to film, script, documentary, cast
  \item to vacuum, vacuum cleaner, to sweep, to scrub, mop
  \item detergent, to hang out, clothes horse, clothes peg, to iron
  \item landlord, tenant, to furnish, pantry, damp
\end{itemize}

\begin{tcolorbox}[breakable,colback=teal!4,colframe=teal!35!black,title={Before you move on}]
Chapter? (\textbf{o capítulo}.) Novelist? (\textbf{o romancista}.) Biography? (\textbf{a biografia}.) To reread? (\textbf{reler}.) Bookmark? (\textbf{o marcador}.) Poetry? (\textbf{a poesia}.) Poet? (\textbf{o poeta}.) Line? (\textbf{o verso}.) Rhyme? (\textbf{a rima}.) To recite? (\textbf{declamar}.) Librarian? (\textbf{o bibliotecário}.) To borrow? (\textbf{requisitar}.) Publisher? (\textbf{a editora}.) Translation? (\textbf{a tradução}.) Translator? (\textbf{o tradutor}.) Episode? (\textbf{o episódio}.) Season? (\textbf{a temporada}.) Soap opera? (\textbf{a telenovela}.) Subtitled? (\textbf{legendado}.) Dubbing? (\textbf{a dobragem}.) (film) director? (\textbf{o realizador}.) To film? (\textbf{filmar}.) Script? (\textbf{o guião}.) Documentary? (\textbf{o documentário}.) Cast? (\textbf{o elenco}.) To vacuum? (\textbf{aspirar}.) Vacuum cleaner? (\textbf{o aspirador}.) To sweep? (\textbf{varrer}.) To scrub? (\textbf{esfregar}.) Mop? (\textbf{a esfregona}.) Detergent? (\textbf{o detergente}.) To hang out? (\textbf{estender}.) Clothes horse? (\textbf{o estendal}.) Clothes peg? (\textbf{a mola}.) To iron? (\textbf{engomar}.) Landlord? (\textbf{o senhorio}.) Tenant? (\textbf{o inquilino}.) To furnish? (\textbf{mobilar}.) Pantry? (\textbf{a despensa}.) Damp? (\textbf{a humidade}.)
\end{tcolorbox}
